\section{Funcionalidades de la plataforma}
\subsection{Acceso y wallet}
La aplicación permite crear una cuenta e iniciar sesión. La wallet se conecta por separado y se vincula al usuario mediante la firma de un mensaje con un código de un solo uso. El servidor comprueba esa firma y registra la dirección del participante.

Las acciones sobre fondos utilizan la wallet vinculada al usuario. La interfaz muestra la dirección, la red y el estado de verificación. Los permisos de gobernanza se consultan en los contratos.

\subsection{Pantallas y operaciones}
\begin{longtable}{L{.23\linewidth}L{.67\linewidth}}
\caption{Funciones de las pantallas de la aplicación}\label{tab:funcionalidades}\\
\toprule Pantalla & Funciones \\ \midrule\endfirsthead
\toprule Pantalla & Funciones \\ \midrule\endhead
Inicio & Consultar generación, consumo, exportación, SUN, obligaciones e ingresos.\\
Hogares & Registrar, consultar y editar ubicación, zona horaria, capacidad solar y batería.\\
Dispositivos & Registrar equipos, seleccionarlos por hogar y consultar lecturas, salud y calidad.\\
Energía & Filtrar por hogar y periodo, revisar gráficas y consultar intervalos.\\
Comunidad & Consultar reglas y nodos, presentar propuestas y solicitar membresía.\\
Créditos SUN & Consultar cantidades, estados de emisión y relación con el intervalo de origen.\\
Liquidación & Revisar lotes, asignaciones, financiación y pruebas; consultar las tareas del operador.\\
Gobernanza & Preparar, programar y ejecutar acciones de la cuenta autorizada.\\
Depósitos y retiros & Depositar ETH, autorizar y depositar SUN, retirar activos disponibles y cobrar ingresos.\\
Compensación & Seleccionar consumo entregado, autorizar SUN y registrar su quema.\\
Transacciones & Consultar la actividad de la sesión y el registro de operaciones del servidor.\\
Configuración & Consultar el perfil, el estado de servicios y la zona horaria utilizada.\\
\bottomrule
\end{longtable}

\subsection{Operaciones del productor}
El productor recibe SUN por la inyección aprobada. Para usarlos en la liquidación, autoriza una cantidad al contrato de custodia y luego ejecuta el depósito. La aplicación muestra los SUN en wallet, los disponibles en depósito y los reservados.

Cuando se ejecuta una asignación, el contrato acredita ETH al productor. Este ingreso aparece como reclamable y se envía a su wallet al firmar el retiro. También puede recuperar los SUN depositados que estén disponibles.

\subsection{Operaciones del consumidor}
El consumidor consulta la energía que tiene asignada y su coste. Deposita ETH con un precio máximo permitido; el contrato reserva los fondos correspondientes a cada asignación. Al liquidarse, recibe SUN.

En Compensación selecciona el consumo entregado, autoriza el uso de los SUN y firma la operación que los quema. El registro resultante vincula esa cantidad con el consumo aprobado. El consumidor también puede retirar ETH disponible de su depósito.

\subsection{Gobernanza}
La cuenta de gobernanza administra las autorizaciones de la comunidad y la finalización de lotes. La aplicación prepara la acción y la wallet firma su programación y ejecución. La finalización confirmada habilita el procesamiento de las emisiones correspondientes.

\subsection{Revisión y firma}
Antes de una operación, la aplicación presenta red, cuenta, contrato, valor y argumentos. La firma se realiza en MetaMask. Autorizar un gasto y transferir los activos son acciones separadas: \code{approve_sun} autoriza el depósito y \code{escrow_sun} lo ejecuta. Para el consumo, \code{approve_retirement} autoriza la quema y \code{retire} la registra.

\subsection{Estados de la operación}
El usuario puede seguir cuatro estados de negocio: intervalo calculado, lote aprobado, crédito emitido y asignación liquidada. El estado \code{FINALIZED} pertenece al cálculo o al lote según la pantalla; \code{ISSUED} identifica un crédito emitido y \code{SETTLED} una asignación liquidada.

La financiación se consulta por separado. El estado \textit{Con fondos} indica que hay recursos suficientes o una reserva constituida. La entrega se confirma mediante el estado de liquidación.

\subsection{Unidades y saldos}
Un SUN representa un kWh de energía elegible aprobada. El token utiliza dieciocho decimales para registrar fracciones. Por ejemplo, una asignación de 0,108 kWh corresponde a 0,108 SUN.

Los saldos de wallet, depósito, reserva e ingreso corresponden a ubicaciones y usos distintos de los activos. La consulta separada permite conocer qué cantidad puede depositarse, retirarse o utilizarse en una asignación.
